\subsection{某一状态附近的线性稳定性}

非线性动力系统中的一个标准做法，是在某个不动点或轨道附近进行线性化分析。
这一做法在本手稿中的用途很克制：
它帮助我们区分，某个局部状态表现为“恢复型”行为，还是表现为“失稳型”行为。

\begin{remark}
对于本书而言，局部线性稳定性之所以有价值，
是因为它使读者能够提出一个更准确的问题：一个小扰动会缩小，还是会放大？
仅仅这一问题，就已经改善了干预设计思路。
如果小扰动会衰减，系统就是鲁棒的；
如果小扰动会增长，系统就处在一个不稳定边缘附近，此时极小变化也可能造成巨大影响。
\end{remark}